\section{Envolvente Compleja para procesos estacionarios}

Para el estudio de procesos estocásticos modulados se ha utilizado extensamente
la envolvente compleja. En este Apéndice se discute el sustento matemático de ese análisis. 


Comenzamos considerando un proceso aleatorio $x(t) + j y(t)$ a valores complejos,
donde  $x(t)$, $y(t)$ son procesos estacionarios de media 0, que cumplen siguientes condiciones: 
\begin{equation*}
    \label{eq:imp}
    \begin{cases}
R_y(\tau)&=\ R_x(\tau);\\
                    R_{xy}(\tau)&=\ E\left[x(t)y(t-\tau)\right]\text{ impar en
                    $\tau$}.
                    \end{cases}
                \tag{\#}
\end{equation*}
La primera condición, ya impuesta en la Sección \ref{sec.modest},  implica que $x(t)$, $y(t)$ tienen
características similares, en particular tienen la misma potencia media. La segunda condición es menos intuitiva, notar que es más general que la hipótesis de incoherencia ($R_{xy}\equiv 0$) considerada en la Sección \ref{sec.modest}. Una versión  equivalente sería  $R_{yx}(\tau)=-R_{xy}(\tau)$, ya que
\[R_{yx}(\tau)=E\left[y(t)x(t-\tau)\right]=E\left[x(t)y(t+\tau)\right]=R_{xy}(-\tau).\]

 Otra forma de representar
\eqref{eq:imp} es en frecuencia:

\[\begin{cases} S_y(f)&= \ S_x(f);\\
                S_{xy}(f)&= \ \mathcal{F}\left[R_{xy}(\tau)\right]\text{ es imaginario
                puro}.
              \end{cases}\tag{\#}\]


Introducimos ahora la modulación, rotando la ``envolvente
compleja'' $x(t)+jy(t)$, con velocidad $\omega_c$. Definimos:
\[u(t)+jv(t)=[x(t)+jy(t)]\e^{j\omega_ct}.\]

\begin{equation*}
    \left[\begin{array}{c}
            u(t)\\
            v(t)
          \end{array}\right]
    = \underbrace{\left[\begin{array}{cc}
                \cos(\omega_ct) & -\sen(\omega_ct)\\
                \sen(\omega_ct) & \cos(\omega_ct)
            \end{array}\right]}_{\text{matriz de rotación}}
    \left[\begin{array}{c}
            x(t)\\
            y(t)
          \end{array}\right].
\end{equation*}

El siguiente Lema generaliza el análisis de procesos modulados de la Sección \ref{sec.modest}.


\begin{lemaSN}
    Si $\left(x(t),y(t)\right)$ son procesos estacionarios, de media 0, que cumplen las propiedades de \eqref{eq:imp},
    entonces $\left(u(t),v(t)\right)$ también  son
    procesos estacionarios de media 0 que cumplen  \eqref{eq:imp}. En particular:
    \begin{align*}
%    \[\left\lbrace\begin{gathered}
                R_u(\tau)=R_v(\tau)&=R_x(\tau)\cos(\omega_c\tau)+R_{xy}(\tau)\sen(\omega_c\tau),\\
                R_{uv}(\tau)&=R_{xy}(\tau)\cos(\omega_c\tau)-R_x(\tau)\sen(\omega_c\tau). %\text{ (impar en $\tau$)}
              \end{align*}%gathered}\right.\tag{\#}\]

        \begin{proof} Es inmediato que
        $u(t)=x(t)\cos(\omega_ct)-y(t)\sen(\omega_ct)$ tiene media
        0 si $x(t)$, $y(t)$ la tienen. Estudiamos ahora la autocorrelación de
        $u(t)$.
            \begin{align*}
                E\left[u(t)u(t-\tau)\right]  = & \ R_x(\tau)\cos(\omega_ct)\cos\left(\omega_c(t-\tau)\right)
                                 + R_y(\tau)\sen(\omega_ct)\sen\left(\omega_c(t-\tau)\right) \\
                                & - R_{xy}(\tau)\cos(\omega_ct)\sen\left(\omega_c(t-\tau)\right)
                                  - R_{yx}(\tau)\sen(\omega_ct)\cos\left(\omega_c(t-\tau)\right)\\
                                 = & \ R_x(\tau)\left[\cos(\omega_ct)\cos\left(\omega_c(t-\tau)\right)+\sen(\omega_ct)\sen\left(\omega_c(t-\tau)\right)\right] +\\
                                & + R_{xy}(\tau)\left[\sen(\omega_ct)\cos\left(\omega_c(t-\tau)\right)
                                - \cos(\omega_ct)\sen\left(\omega_c(t-\tau)\right)\right]    \\
                                 = & \
                                 R_x(\tau)\cos(\omega_c\tau)+R_{xy}(\tau)\sen(\omega_c\tau).
            \end{align*}
            Aquí la primera igualdad es por definición, la segunda utiliza  \eqref{eq:imp}, y la
            tercera identidades trigonométricas. Como el resultado depende solo de $\tau$, el proceso $u(t)$ es
            estacionario en sentido amplio. El análisis para
            $v(t)$, y para la correlación cruzada $R_{uv}$ es
            análogo.
        \end{proof}
\end{lemaSN}

\subsection*{Aplicación a la envolvente compleja}
Recordemos el procedimiento para hallar la
envolvente compleja $x_{ce}(t)$ para señales pasabanda
determinísticas $x_c(t)$. En Fourier,
%$\hat{X}_{ce}(f)=2u(f+f_c)\hat{X}_c(f+f_c)$:
truncábamos a frecuencias positivas, $2 u(f)\hat{X}_c(f)$, y
trasladamos hacia atrás $f_c$ en frecuencia. Un diagrama de
bloques es:
\begin{figure}[htb]
    \centering\begin{minipage}{\textwidth}
    \bpgf{Imagenes/ruidoPasa/ruidoPasa-2}
    \begin{blockdiagram}
        \matrix[row sep=5mm, column sep=10mm]{
        \node [etiqueta] (inicio) {$x_c(t)\quad $}; &  \node[block] (filtro) {$2u(f)$}; & \node [operator] (mult) {$\times$}; & \node[etiqueta] (salida) {$\quad x_{ce}(t)$};\\
        & & \node[etiqueta] (mod) {$\mepi$};&\\
        };
        \draw [->] (inicio) -- (filtro);
        \draw [->] (filtro) -- (mult);
        \draw [->] (mult) -- (salida);
        \draw [->] (mod) -- (mult);
    \end{blockdiagram}
    \endpgfgraphicnamed\end{minipage}
\end{figure}
\FloatBarrier

También recordar que $2u(f)=1+j\hat{H}(f)$, siendo
$\hat{H}(f)=-j\sgn(f)$ el filtro de Hilbert:

\begin{figure}[htb]
    \centering \begin{minipage}{\textwidth}
    \bpgf{Imagenes/ruidoPasa/ruidoPasa-3}
    \begin{blockdiagram}
        \matrix[row sep=5mm, column sep=10mm]{
        \node [etiqueta] (inicio) {$x_c(t)\quad $}; &  \node [punto] (p) {}; & & & \node[operator] (sum) {$+$};& \node [operator] (mult) {$\times$}; & \node[etiqueta] (salida) {$\quad x_{ce}(t)$};\\
        & & \node[block] (hil) {$\mathcal{H}$}; & \node[block] (j) {$j$}; & & \node[etiqueta] (mod) {$\mepi$};&\\
        };
        \draw [->] (inicio) -- (p) -- (sum);
        \draw [->] (p) |- (hil);
        \draw [->] (hil) -- (j);
        \draw [->] (j) -| (sum);
        \draw [->] (sum) -- (mult);
        \draw [->] (mult) -- (salida);
        \draw [->] (mod) -- (mult);
    \end{blockdiagram}
    \endpgfgraphicnamed\end{minipage}
\end{figure}
\FloatBarrier

Aplicamos ahora el diagrama de la figura anterior al ruido pasabanda. O sea, sustituimos la
entrada $x_c(t)$ por un ruido $n(t)$ que cumple $S_n(f)=0$ fuera de la
banda $[f_c-B,f_c+B]$. Tenemos: 
\[
n_{ce}(t)=n_i(t)+jn_q(t) = \left(n + j\mathcal{H}n\right)
\e^{-j\omega_ct},\]
\begin{equation*}
    \left[\begin{array}{c}
            n_i(t)\\
            n_q(t)
          \end{array}\right]
    = \left[\begin{array}{cc}
                \cos(\omega_ct) & \sen(\omega_ct)\\
                -\sen(\omega_ct) & \cos(\omega_ct)
            \end{array}\right]
    \left[\begin{array}{c}
            n(t)\\
            \mathcal{H}n(t)
          \end{array}\right].
\end{equation*}

\paragraph{Notas:} Sea $\check{n}=\mathcal{H}n$ (transformada de Hilbert del ruido).
Como $\mathcal{H}$ es LTI, $\Rightarrow \check{n}(t)$ es un
proceso estacionario, y se cumple:
\[\begin{cases} S_{\check{n}}(f)& =\left|\hat{H}(f)\right|^2S_n(f)=S_n(f);\\
                S_{n\check{n}} &= \underbrace{\overline{\hat{H}(f)}}_{\text{\tiny imag puro}}\underbrace{S_n(f)}_{\text{\tiny  real}}    
    \text{  imaginario
                puro}.
              \end{cases}\tag{\#}\]


%\newpage
Por lo tanto $n, \check{n}$ cumplen las hipótesis \eqref{eq:imp} del
Lema. Entonces (cambiando $\omega_c$ por $-\omega_c$), tenemos que
$(n_i,n_q)$ son procesos estacionarios, y
\begin{align*}
    R_{n_i}(\tau)=R_{n_q}(\tau) & = R_n(\tau)\cos(\omega_c\tau)+R_{n\check{n}}\sen(-\omega_c\tau)\\
                    & = R_{n}(\tau)\cos(\omega_c\tau) +
                    R_{\check{n}n}\sen(\omega_c\tau).
\end{align*}
¿Cómo es el espectro de $n_i$, $n_q$?
\begin{align*}
    S_{n_i}(f)  & = S_n(f) \ast \left[\frac{\delta(f-f_c)+\delta(f+f_c)}{2}\right] +
              \left[\hat{H}(f)S_n(f)\right]\ast j\left[\frac{-\delta(f-f_c)+\delta(f+f_c)}{2}\right]\\
            & = \left(S_n(f)\underbrace{\left(\frac{1-j\hat{H}(f)}{2}\right)}_{u(-f)}\right)\ast\delta(f-f_c)+
            \left(S_n(f)\underbrace{\left(\frac{1+j\hat{H}(f)}{2}\right)}_{u(f)}\right)\ast\delta(f+f_c).\\
%\Longrightarrow 
& = \underbrace{S_n(f-f_c)u\left(-(f-f_c)\right)}_{\oldstylenums{1}}+
    \underbrace{S_n(f+f_c)u(f+f_c)}_{\oldstylenums{2}}.
\end{align*}

\noindent Por otra parte, utilizando la segunda condición del lema (cambiando
$\omega_c$ por $-\omega_c$) tenemos
\begin{align*}
    R_{n_in_q} & =R_n(\tau)\sen(\omega_c\tau)-R_{\check{n}n}(\tau)\cos(\omega_c\tau)\\
    \Rightarrow %\begin{aligned}
                    S_{n_in_q}  & = S_n(f) \ast \left[\frac{-j\delta(f-f_c)+j\delta(f+f_c)}{2}\right] - \left[\hat{H}(f)S_n(f)\right] \ast \left[\frac{\delta(f-f_c)+\delta(f+f_c)}{2}\right]\\
                            & = j \left\lbrace\left[S_n(f)\left(\frac{1+j\hat{H}(f)}{2}\right)\right] \ast \delta(f+f_c) - \left[S_n(f)\left(\frac{1-j\hat{H}(f)}{2}\right)\right] \ast \delta(f-f_c)\right\rbrace\\
                            & = j \left\lbrace\underbrace{S_n(f+f_c)u(f+f_c)}_{\oldstylenums{2}}-
                            \underbrace{S_n(f-f_c)u\left(-(f-f_c)\right)}_{\oldstylenums{1}}\right\rbrace.
                %\end{aligned}
\end{align*}
Aparece aquí la {\em diferencia} de los dos espectros
considerados anteriormente. Como consecuencia: si $S_n(f)$ es simétrica respecto de $f_c$, 
$S_{n_in_q}=0 \Rightarrow n_i(t), n_q(t)$ no correlacionados.

Hemos probado entonces el Teorema enunciado en la Sección \ref{sec.ruidopasa}. 

\paragraph{Notas:} 
\begin{itemize}
\item  $n_i$, $n_q$ son estacionarios, pasabajos, con
$S_{n_i}(f)=S_{n_q}(f)$ simétrico respecto de $f=0$. Cuando son no
correlacionadas entre sí, corresponde a un ruido pasabanda 
simétrico respecto de $f_c$. Un ruido pasabanda asimétrico indica correlación
(informalmente, dependencia) entre las componentes en fase y
cuadratura. En un extremo de dependencia, cuando
$n_q=\mathcal{H}{n_i}$, el ruido pasabanda será SSB (por encima de
$f_c$).%
\item De acuerdo al Lema, $n_i$, $n_q$
``heredan" la propiedad \eqref{eq:imp} que cumplían $n$,
$\check{n}$. Se deduce de aquí que si uno realiza una nueva modulación del tipo 
\[n'(t)=n_i(t)\cos(\omega'_ct)-n_q(t)\sen(\omega'_ct)\]
se obtendrá un proceso estacionario pasabanda, en banda
$|f|\in[f'_c-B,f'_c+B]$. 
\end{itemize}

%


